\documentclass[11pt]{article}
\usepackage[T1]{fontenc}
\usepackage{lmodern}
\usepackage{amsmath,amssymb,amsthm}
\usepackage{graphicx}
\usepackage[margin=0.88in]{geometry}
\usepackage[colorlinks=true,linkcolor=blue,citecolor=blue,urlcolor=blue]{hyperref}

\newtheorem{theorem}{Theorem}
\newtheorem{corollary}{Corollary}
\newtheorem{remark}{Remark}
\newcommand{\R}{\mathbb R}
\newcommand{\sphere}{\mathbb S}
\newcommand{\Harm}{\mathcal H}
\newcommand{\Poly}{\mathcal P}

\title{Finite Families of Subspaces of Codimension at Least Two\\
Are Never Helmholtz Uniqueness Sets}
\author{Run \texttt{fa\_banach\_001}, agent \texttt{agent\_lane\_11}}
\date{17 August 2026}

\begin{document}
\maketitle

\section*{Source Question and Status}

Fern\'andez-Bertolin, Gr\"ochenig, and Jaming prove nonuniqueness for an
arbitrary finite collection of lines in $\R^d$.  Immediately afterwards
they write: ``We do not know what happens for restrictions of the Helmholtz
equation to $k$-dimensional subspaces when $2\leq k\leq d-2$ and $d\geq4$''
\cite[p.~18]{FGJ}.

\begin{center}
\includegraphics[width=0.96\linewidth]{figures/open_question_crop.png}
\end{center}

In the fixed finite-family setting of their preceding general problem, the
answer is always negative.  In fact, no finite family of linear subspaces
of codimension at least two is a uniqueness set, for any real wave number.
The same construction disproves the associated sphere/subspace Heisenberg
uniqueness property.  Status: \emph{candidate full counterexample theorem,
likely valid}, subject to human review.

\section*{Proof Intuition}

Vanishing on a whole $k$-plane imposes more than finitely many point
conditions, but at a fixed homogeneous degree it still imposes only
$O(m^{k-1})$ linear conditions.  Degree-$m$ spherical harmonics in $d$
variables have dimension $O(m^{d-2})$.  When $k\leq d-2$, a finite number
of restriction maps therefore cannot detect every high-degree spherical
harmonic.  A common kernel harmonic gives a harmonic counterexample
directly.  Multiplying it by the regular radial Helmholtz power series gives
the same counterexample at every wave number, without a singularity at the
origin.

\section*{Main Result}

Let $\Harm_m^d$ be the real vector space of homogeneous harmonic
polynomials of degree $m$ on $\R^d$, and let $\Poly_m(E)$ be the homogeneous
degree-$m$ polynomials on a linear space $E$.

\begin{theorem}[finite subspaces are not unique]\label{thm:main}
Let $d\geq3$, let $E_1,\ldots,E_N\subset\R^d$ be linear subspaces with
$1\leq k_j:=\dim E_j\leq d-2$, and let $\lambda\in\R$.  There exists a
nonzero real entire function $u$ on $\R^d$ such that
\[
 (\Delta+\lambda^2)u=0,
 \qquad
 u|_{E_1\cup\cdots\cup E_N}=0.
\]
Consequently, for $d\geq4$ and $2\leq k\leq d-2$, no finite collection of
$k$-dimensional linear subspaces is a uniqueness set for Helmholtz
solutions.
\end{theorem}

\begin{proof}
For each $m\geq1$, consider the direct-sum restriction map
\[
 \mathcal R_m:\Harm_m^d\longrightarrow
       \bigoplus_{j=1}^N\Poly_m(E_j),
 \qquad H\longmapsto(H|_{E_1},\ldots,H|_{E_N}).
\]
The Fischer decomposition of homogeneous polynomials gives
\begin{align}
 \dim\Harm_m^d
 &=\binom{m+d-1}{d-1}-\binom{m+d-3}{d-1}\label{eq:hdim}\\
 &=\binom{m+d-2}{d-2}+\binom{m+d-3}{d-2}
   =\frac{2}{(d-2)!}m^{d-2}+O(m^{d-3}).\nonumber
\end{align}
On the other hand,
\begin{equation}\label{eq:rdim}
 \dim\Poly_m(E_j)=\binom{m+k_j-1}{k_j-1}=O(m^{k_j-1})
 =O(m^{d-3}).
\end{equation}
Since $N$ is fixed, equations \eqref{eq:hdim}--\eqref{eq:rdim} show that
for all sufficiently large $m$,
\[
 \dim\Harm_m^d>
 \sum_{j=1}^N\dim\Poly_m(E_j).
\]
Hence $\ker\mathcal R_m$ contains a nonzero real polynomial $H$.  Thus
$H|_{E_j}=0$ for every $j$.

Put $s=|x|^2$ and define the entire scalar function
\begin{equation}\label{eq:radial}
 \Phi_m(s)=\sum_{\ell=0}^{\infty}
 \frac{(-\lambda^2s/4)^\ell}{\ell!\,(m+d/2)_\ell},
 \qquad u(x)=H(x)\Phi_m(|x|^2),
\end{equation}
where $(a)_\ell$ is the rising factorial.  The series has infinite radius
of convergence and $\Phi_m(0)=1$, so $u$ is entire and nonzero.  Since $H$
is homogeneous of degree $m$ and harmonic,
\begin{equation}\label{eq:laplace}
 \Delta\bigl(H(x)\Phi(|x|^2)\bigr)
 =H(x)\{4s\Phi''(s)+(4m+2d)\Phi'(s)\}.
\end{equation}
If $a_\ell$ denotes the coefficient of $s^\ell$ in \eqref{eq:radial},
then
\[
 4(\ell+1)(\ell+m+d/2)a_{\ell+1}+\lambda^2a_\ell=0.
\]
Substitution in \eqref{eq:laplace} gives
$(\Delta+\lambda^2)u=0$.  Finally, the factor $H$ shows that $u$ vanishes
identically on every $E_j$.
\end{proof}

\section*{Heisenberg Uniqueness Consequence}

\begin{corollary}\label{cor:hup}
Under the assumptions of Theorem~\ref{thm:main},
$(\sphere^{d-1},E_1\cup\cdots\cup E_N)$ is not a Heisenberg uniqueness
pair.  More strongly, there are two distinct finite positive measures on
$\sphere^{d-1}$ whose Fourier transforms agree on every $E_j$.
\end{corollary}

\begin{proof}
Choose $m\geq1$ and $0\neq H\in\ker\mathcal R_m$ as above, and let
$Y=H|_{\sphere^{d-1}}$.  After scaling, the measures
\[
 d\mu_\pm(\theta)=
 \left(1\pm\frac{Y(\theta)}{\|Y\|_\infty}\right)d\sigma(\theta)
\]
are positive and distinct.  The Funk--Hecke formula gives, for a nonzero
constant $c_{m,d}$ and $r>0$,
\[
 \widehat{Y\,\sigma}(r\omega)
 =c_{m,d}\,r^{-(d-2)/2}J_{m+(d-2)/2}(r)Y(\omega).
\]
If $r\omega\in E_j$, then $Y(\omega)=0$.  At the origin the transform also
vanishes because every positive-degree spherical harmonic has mean zero.
Thus $\widehat{\mu_+}=\widehat{\mu_-}$ on each $E_j$.
\end{proof}

\section*{Scope, Novelty Check, and Verification}

The result answers the source's stated intermediate-dimensional linear-
subspace case for every finite configuration.  It does not claim
nonuniqueness for infinitely many subspaces or for arbitrary curved
lower-dimensional manifolds.  The codimension threshold is sharp for this
argument: when $\dim E=d-1$, the source proves uniqueness for two generic
intersecting hypersurfaces, and the domain and restriction dimensions have
the same polynomial order.

The run's cheap indexes contained no answer.  Exact-sentence and concept
searches through 17 August 2026 found the source, its published version,
and related HUP/PDE literature, but no later answer to this finite-family
question.  The companion script checks the exact dimension identities,
finds kernel-producing degrees across a grid of dimensions and family
sizes, and verifies the radial coefficient recurrence using rational
arithmetic.  These are consistency checks; the proof above is exact.

\section*{Human Review Recommendation}

Check that the source's open sentence inherits the finite-family setting
from the preceding general problem, and verify the rank bound
$\operatorname{rank}(H\mapsto H|_E)\leq\dim\Poly_m(E)$.  Once those two
points are accepted, the dimension gap and the radial recurrence settle the
linear-subspace problem.

\begin{thebibliography}{9}
\small
\bibitem{FGJ}
A. Fern\'andez-Bertolin, K. Gr\"ochenig, and P. Jaming,
\emph{From Heisenberg uniqueness pairs to properties of the Helmholtz and
Laplace equations}, Journal of Mathematical Analysis and Applications
\textbf{469} (2019), 202--219; arXiv:1711.05520;
doi:10.1016/j.jmaa.2018.09.008.
\end{thebibliography}

\end{document}
